\documentclass[11pt,letterpaper]{article}
\usepackage[margin=.79in,top=.72in,bottom=.72in]{geometry}
\usepackage{lmodern,microtype}
\usepackage{amsmath,amssymb,amsthm}
\usepackage[hidelinks]{hyperref}
\setlength{\parindent}{0pt}
\setlength{\parskip}{4pt}
\newtheorem{theorem}{Theorem}
\hypersetup{pdftitle={Search Under Disorder and Query Entropy},pdfauthor={Jonathan R. Landers}}
\begin{document}
\begin{center}
{\Large\bfseries Search Under Disorder and Query Entropy}\\[2pt]
{\small A two-budget note on repeated comparison search}\\[5pt]
Jonathan R. Landers \qquad October 8, 2026
\end{center}

The disorder potential in \emph{Disorder Potential and the Cost of Search}
limits how far entries can drift from their sorted ranks. It says nothing
about which targets will be requested. Even a sorted array needs logarithmic
worst-case search when every key or gap is possible. This note gives the
targets a known distribution by rank. The resulting bound separates two
costs: disorder controls the comparisons needed to build a sorted index,
while target uncertainty controls the expected comparisons after it exists.

\paragraph{Array promise and query prior.}
Let $A=(a_1,\ldots,a_n)$ contain distinct real keys, and let $x_i$ be
the sorted rank of $a_i$. As in the paper, assume the supplied promise
\[
 V(A)=\frac12\sum_{i=1}^n(x_i-i)^2\le E,\qquad
 D_E=\min\left\{n-1,
 \left\lfloor\frac{\sqrt{1+8E}-1}{2}\right\rfloor\right\},
 \qquad R_E=D_E+1.
\]
The displacement lemma gives $|x_i-i|\le D_E$: if one entry moves by
$d$, the other displacements sum to $-d$, so $2V(A)\ge d(d+1)$.
Without preprocessing, one target costs
$O(\log_2(n+1)+R_E)$ comparisons in the worst case.

A target lies in one of the $2n+1$ ordered cells
\[
 g_0,\ k_1,\ g_1,\ k_2,\ldots,k_n,\ g_n,
\]
where $k_j$ is the key of rank $j$ and $g_j$ is the gap after it.
Let $w$ be known probabilities on these cells, with
$H(w)=-\sum_c w_c\log_2 w_c$ and $0\log_2 0=0$.
Conditional on the fixed array, targets are independent draws with
this cell law. Their values within a gap do not affect comparisons.
The law concerns \emph{target ranks}, not the numerical distribution of
array values.

Preprocessing may compare array keys and retain an $O(n)$-word index.
During a query, only target-to-array-key comparisons reveal order; they are
charged, and the algorithm must answer correctly for every real target,
including cells of probability zero. Define $\mathcal T_q(n,E,w)$ as the
minimum, over these algorithms, of the worst-case over $V(A)\le E$ of
preprocessing comparisons plus expected comparisons for $q$ targets.

\begin{theorem}[Two-budget search bound]
For $n\ge2$ and $q\ge1$, with an absolute constant $C$,
\[
 \boxed{\displaystyle
 \frac{qH(w)}{\log_2 3}
 \ \le\ \mathcal T_q(n,E,w)\
 \le\ C\min\left\{
 q\bigl(\log_2(n+1)+R_E\bigr),\
 n\log_2(R_E+1)+q\bigl(H(w)+2\bigr)
 \right\}.}
 \tag{1}
\]
When $E=0$, the sorted-index preprocessing in the second term uses
zero comparisons.
\end{theorem}

If $w$ is spread roughly uniformly, $H(w)=\Theta(\log n)$ and the
new term gives no asymptotic saving per target. If queries are
supported on $m$ cells, then $H(w)\le\log_2 m$; after the index is
built, their expected search cost is $O(\log m+1)$, even though the
algorithm remains correct on every target. The index cost is paid once
and can be amortized across many queries.

The saving does not require impossible queries. Suppose $1-\varepsilon$
of the mass is uniform on $m$ favored key cells and $\varepsilon$ is
uniform on all other cells. Then every cell has positive probability,
yet
\[
 H(w)=h_2(\varepsilon)+(1-\varepsilon)\log_2 m
       +\varepsilon\log_2(2n+1-m),
\]
where $h_2$ is binary entropy. With fixed $m$ and
$\varepsilon=1/\log_2 n$, this is $O(1)$: after indexing,
the expected comparisons per target are constant, while rare targets
still receive correct answers.

\newpage
\paragraph{Proof of the upper bound.}
The first choice in (1) repeats the potential-based search: midpoint
comparisons shrink the candidate interval up to a margin $2D_E$,
followed by a scan of $O(R_E)$ positions.

For the second choice, build a sorted index of the array positions.
Since the entry of rank $j$ occurs by position $j+D_E$, begin with a
min-heap containing positions $1,\ldots,\min\{n,R_E\}$. Extract its
minimum, insert the next unread position, and repeat. Before the $j$th
extraction, the heap contains every unread position through
$\min\{n,j+D_E\}$; in particular it contains rank $j$.
Each extraction therefore yields the next key in sorted order.
The heap has at most $R_E$ entries, so the index costs
$O(n\log_2(R_E+1))$ array-key comparisons and $O(n)$ words.
For $D_E=0$, the array itself is already the index.

Now search this index with an ordered decision tree chosen for $w$.
For every cell $c$ with $w_c>0$, place a dyadic interval of length
$2^{-\ell_c}$ inside its cumulative-probability interval, where
$\ell_c=\lceil\log_2(1/w_c)\rceil+1$. These intervals give an
alphabetic prefix code: their left-to-right order is the cell order,
and their mean depth is less than $H(w)+2$.
Zero-weight cells can be included by an arbitrarily small perturbation.
Every split between consecutive cells is a threshold at some indexed
key: a cut just before or after rank $j$ uses $y<k_j$ or
$y\le k_j$, respectively. One target-to-key comparison
therefore implements the binary branch. Thus the mean search cost is $O(H(w)+2)$, proving the
second choice. This is the classical entropy-sensitive search-tree
construction.\footnote{K.~Mehlhorn, \emph{Nearly Optimal Binary Search
Trees}, Acta Informatica 5 (1975), 287--295.
\href{https://people.mpi-inf.mpg.de/~mehlhorn/ftp/mehlhorn3.pdf}{Author-hosted paper}.}

\paragraph{Why the entropy term is necessary.}
Fix any array and any preprocessing transcript. A comparison of the
target with an array key has at most three outcomes. A correct search
tree must distinguish its possible key and gap cells: an absent leaf
cannot cover gaps on both sides of a key without also admitting that
key. The leaf depths therefore form a ternary prefix code for $w$,
whose expected depth is at least $H(w)/\log_2 3$. Independence makes
the same conditional bound hold after any history of earlier queries.
Summing over $q$ proves the left side of (1).

\paragraph{When the saving pays.}
Writing $B_E=O(n\log_2(R_E+1))$, the two strategies give an average
cost per target of at most
\[
 O\!\left(\min\left\{\log_2(n+1)+R_E,\;
                  \frac{B_E}{q}+H(w)+1\right\}\right).
 \tag{2}
\]
For $E=0$, $B_E=0$ and a known query prior helps immediately.
For a disordered array, enough queries must share the index to repay
its construction. This is an amortized gain in expected search depth,
not a claim that the worst-case target has become easier.

\paragraph{What the bound changes.}
The original $\Omega(\log n)$ obstruction concerns a worst-case target;
it remains when all cells are possible and no query prior is used.
A law for the \emph{array's} potential, beta-shaped or otherwise, can
reduce the displacement work but does not lower target entropy.
Here the search-depth saving is an expectation over a stated query law,
and the sorted index that makes it possible is charged explicitly.
A sharper no-preprocessing result of the form
$O(H(w)+R_E)$ per query would require a separate search argument;
(1) makes no such claim.

\paragraph{Takeaway.}
Disorder sets the one-time price of learning the array's order, and
query entropy sets the average price of using that order. When targets
concentrate on a few ranks and enough queries share the index, the
recurring $\log n$ search term becomes $H(w)+O(1)$ without giving up
correct answers for rare targets.

\smallskip
{\footnotesize\emph{Context.} The potential, displacement lemma, and
single-search theorem are from Landers,
\emph{Disorder Potential and the Cost of Search}, version 2.
The heap-based index construction comes from Landers,
\emph{Search Under a Disorder Budget}.}
\end{document}
